\section{Experiments}
\label{sec:experiments}

\paragraph{Research questions.}
Organize the evaluation around explicit questions, such as effectiveness, efficiency, robustness, and generalization.

\paragraph{Setup.}
Document datasets, model versions, hardware, software, random seeds, hyperparameters, stopping criteria, and statistical procedures. Store machine-readable experiment configurations outside the paper when possible and cite the corresponding artifact.

\begin{table}[t]
  \centering
  \caption{Example results. Higher quality is better; lower cost is better. Report values as mean $\pm$ standard deviation over independent runs.}
  \label{tab:example-results}
  \begin{tabular}{lcc}
    \toprule
    Method & Quality $\uparrow$ & Cost $\downarrow$ \\
    \midrule
    Baseline A & $71.2 \pm 0.8$ & $1.00\times$ \\
    Baseline B & $73.5 \pm 0.6$ & $1.18\times$ \\
    \methodname & $\mathbf{76.4 \pm 0.5}$ & $0.91\times$ \\
    \bottomrule
  \end{tabular}
\end{table}


\paragraph{Main results.}
\Cref{tab:example-results} is a placeholder for the primary comparison. Report uncertainty, the number of runs, and practical effect sizes in addition to point estimates.

\begin{figure}[t]
  \centering
  \fbox{\rule{0pt}{1.25in}\rule{0.9\columnwidth}{0pt}}
  \caption{Replace this box with a vector graphic or a high-resolution raster image. The caption should state the takeaway without requiring readers to search the main text.}
  \label{fig:placeholder}
\end{figure}

\paragraph{Ablations and failure analysis.}
Evaluate which components matter and characterize cases where the method does not work. Negative results and boundary conditions often strengthen the paper's claims.
